\chapter{Revisión de la Literatura}

\noindent La pregunta sobre si las cámaras de velocidad reducen realmente la siniestralidad vial ha sido abordada en distintos contextos y con metodologías variadas. Antes de presentar la estrategia empírica de esta tesis, es útil situarla dentro de esa literatura: entender qué se ha encontrado en otros países cuando se implementan sistemas de fotodetección, cómo se ha definido el área de influencia y la duración de sus efectos, qué recomendaciones existen sobre su diseño institucional, y qué se ha estudiado ya específicamente en la Ciudad de México. Esta sección recorre primero la evidencia internacional sobre cámaras de velocidad y control automatizado, después revisa los trabajos que discuten su ámbito espacial y temporal de impacto, continúa con las guías de implementación de estos programas y, finalmente, resume los estudios previos para la capital mexicana, destacando las lecciones que retomo y las brechas que motivan el enfoque de esta investigación.

\section{Evidencia internacional sobre cámaras de velocidad}

\noindent La literatura internacional sobre cámaras de velocidad y control automatizado es relativamente consistente: cuando estos dispositivos se diseñan e implementan con buen método y en los lugares adecuados, tienden a reducir los incidentes viales, en particular aquellos con personas lesionadas. Buena parte de la evidencia más sólida proviene del Reino Unido, donde se cuenta con registros detallados de siniestros (STATS19) y con programas de instalación de cámaras suficientemente grandes como para aplicar métodos cuasi–experimentales.\\

Li et al. (2013) analizan 771 cámaras de velocidad utilizando Propensity Score Matching (PSM). El punto de partida es que las cámaras no se ubican al azar, sino en sitios con alto historial de incidentes y ciertas características viales; por ello, comparan cada sitio tratado con controles similares en términos de choques previos, tipo de vía, límite de velocidad, longitud del tramo, afluencia promedio diaria (AADT) y número de intersecciones menores. Una vez corregida la selección, encuentran reducciones claras en incidentes con lesiones personales (PICs) en la vecindad inmediata de los dispositivos: alrededor de -27.5\% a 200 metros y -18.5\% a 1 kilómetro. El estudio además introduce dos ideas que serán centrales en trabajos posteriores: la posible "migración" de accidentes (por desvío de rutas o cambios de conducción) y el llamado efecto \"canguro", que describe el patrón de frenar abruptamente antes de la cámara y acelerar de nuevo después.\\ 

% Para el efecto canguro hay que agregar referencias de Elvik 1997 y Thomas et al. 2008

Graham et al. (2019) estiman el efecto causal de las cámaras en Inglaterra utilizando un enfoque doblemente robusto dentro de un enfoque estadístico bayesiano, combinando modelos de Propensity Score y Outcome Regression (OR) para controlar por sesgos de selección y variables no observadas asociados a proceso de selección no aleatorio de las cámaras. Sus resultados muestran que, en tramos de entre 400 y 1500 metros, los radares reducen en promedio 15\% de las colisiones con lesiones personales. No hacen pruebas en diferentes longitudes de tramos, y la longitud es una de las variables para el emparejamiento. Resaltan también que cuando no se controla por ninguna variable el efecto es del 18\%, y del 34\% cuando no se hace emparejan las unidades.\\ 

Li et al. regresan en 2020 a estudiar cómo afectó el tiempo a  los resultados. En este caso combinan PSM con un diseño de Difference-in-Differences (DID) y siguen los sitios entre 1999 y 2016, es decir, tres años previos a la instalación y 12 años posteriores. Dividen el periodo post-tratamiento en cuatro ventanas temporales para ver si el impacto se atenúa. Para los PICs, el patrón muestra que las reducciones arrancan en torno a -26.8\% en los primeros años posteriores a la instalación y se van atenuando hasta -12.2\% en el perido más tardío de su análisis. Esto sugiere que, si no se acompaña de gestión y un continuo esfuerzo por la efectividad de las políticas, el efecto de las cámaras tiende a disiparse con el tiempo. En accidentes fatales y graves (FSCs) no encuentran tendencias robustas, en buena medida por la baja incidencia y la dispersión de los eventos.\\ 

Mountain et al. (2005) ofrecen un punto de comparación útil al analizar no solo cámaras de velocidad, sino un conjunto amplio de esquemas de gestión de velocidad en vías urbanas con límite de 30 mph. Su base de datos incluye 79 esquemas con cámaras fijas y 71 intervenciones de ingeniería (31 con deflexiones horizontales y 39 con dispositivos verticales, como topes y "cushions"). Mediante un enfoque de before–after con corrección Empirical Bayes, los autores ajustan por tendencias generales en siniestros, por regresión a la media y por "migration", es decir, el posible desplazamiento de accidentes hacia tramos vecinos del corredor. Encuentran que todas las intervenciones reducen las lesiones personales, pero con magnitudes distintas: los esquemas verticales son los más efectivos, con caídas cercanas al 44\% en personal injuries por kilómetro y año, aproximadamente el doble del efecto asociado a las cámaras de velocidad, que ronda el 22\%. Las medidas horizontales son las menos efectivas. Al mismo tiempo, Mountain et al. muestran que las cámaras son particularmente eficaces para reducir el porcentaje de conductores que exceden el límite y la variabilidad de velocidades, sin provocar reducciones tan abruptas como las que generan los dispositivos verticales. Esto las vuelve especialmente adecuadas en contextos donde se busca moderar la velocidad y la severidad de los siniestros, pero donde intervenciones muy intrusivas ---como topes o cushions--- no son viables o deseables por sus costos operativos y de movilidad.\\ 

\subsection{Rangos de efectividad}
\noindent Un punto central es que las cámaras de velocidad tienen un rango de acción espacial acotado. No son instrumentos que cambien de forma homogénea los patrones de siniestralidad en toda la ciudad, sino intervenciones muy locales cuyo efecto se concentra en unos cuantos cientos de metros alrededor del dispositivo. Li et al. (2013) muestran que, en Inglaterra, el impacto de 771 cámaras fijas es claramente mayor en bandas de 200 metros antes y después del punto de instalación; cuando amplían el tramo de análisis a 500 metros y 1 kilómetro por lado, el efecto se mantiene en la misma dirección pero se va diluyendo. Christie et al. (2003), en un estudio con radares móviles, encuentran un patrón similar: usando círculos concéntricos, las caídas en choques con lesionados son muy grandes en los primeros 0-100 metros, más moderadas entre 100-300 metros, y prácticamente inexistentes más allá de 500 metros. En Jordania, Al‑Masaeid et al. (2020) reportan explícitamente que el rango operativo de sus cámaras es de alrededor de 50 metros, lo que ilustra hasta qué punto el efecto esperado es puntual.

\indent Este carácter local también se confirma cuando la unidad de análisis no es un tramo lineal sino un área. Li et al. (2020) estudian "dosis" de tratamiento (una, dos o tres o más cámaras) en discos de radio entre 200 y 1,000 metros alrededor de los dispositivos. Sus resultados muestran que casi todo el beneficio se concentra en radios pequeños: con una sola cámara, la reducción absoluta en accidentes con lesionados por kilómetro y año es mucho mayor a 200 metros que a 600 o 1,000 metros; con dos cámaras, los efectos adicionales solo se observan de forma clara en 200 y 300 metros. Más allá de estos radios, la diferencia entre tener una o varias cámaras prácticamente desaparece. La revisión de De Ceunynck (2017) llega a conclusiones en la misma línea: los efectos de las cámaras fijas son “muy locales” y no se observa evidencia sólida de spillovers hacia sitios no tratados; cuando se detectan efectos a varios kilómetros, suele ser en sistemas de “section control” (control por tramo) que, por diseño, actúan sobre segmentos de varios kilómetros y no sobre puntos aislados. Un sistema de control por tramo es aquel que mide la velocidad promedio de tránsito entre dos puntos, y es especialmente útil en vialidades con accesos controlados de entre uno y 100 kilómetros, pues promueve que se sostenga una velocidad constante a lo largo de todo el trayecto (Job et al. 2020).

\subsection{Ventanas temporales}
\noindent En el plano temporal, los estudios tampoco tratan a las cámaras como intervenciones permanentes cuyo efecto sea constante para siempre, sino que definen ventanas de análisis finitas y comparables. Li et al. (2013) fijan un periodo 1999–2007 precisamente para asegurar, para cada sitio, tres años de información antes y tres después de la instalación del dispositivo, y dentro de ese marco estiman efectos en diferentes bandas espaciales (200, 500 y 1,000 metros). En un trabajo posterior, Li et al. (2020) se concentran en la dinámica del efecto en el tiempo: usan el mismo universo de cámaras, definen tres años previos al tratamiento y dividen los doce años posteriores en cuatro ventanas de igual longitud (2005–2007, 2008–2010, 2011–2013 y 2014–2016). Esta decisión metodológica es deliberada: los autores discuten que ventanas demasiado largas pueden “promediar” periodos muy distintos y ocultar variaciones relevantes, mientras que ventanas muy cortas introducen demasiado ruido aleatorio. Algo similar ocurre en Christie et al. (2003), que trabajan con una serie 1996–2000 y definen periodos before–after de varios años para estabilizar las tasas locales de choques en cada banda espacial.

\indent  Li et al. (2013) sintetizan esta variación en una tabla comparativa: hay trabajos que utilizan periodos muy extensos, de hasta 10 años de datos de choques (por ejemplo, Hess y Polak, 2003; Newstead y Cameron, 2003), típicamente con el objetivo de modelar con más precisión las tendencias de fondo y la regresión a la media. En el extremo opuesto, algunos estudios se basan en horizontes muy cortos, de alrededor de año y medio en total, como Shin et al. (2009) en una autopista urbana de 6.5 millas, lo que permite captar cambios inmediatos pero deja más ruido aleatorio y menos capacidad para separar el efecto de la cámara de otras fluctuaciones del sistema vial. Entre ambos extremos, varios trabajos que intentan ir más allá de un before–after simple convergen en ventanas de orden tres años antes y tres años después de la intervención. Mountain et al. (2004, 2005) y Gains et al. (2004, 2005) adoptan este esquema con datos del Reino Unido, igual que el propio Li et al. (2013), combinándolo con grupos de comparación y métodos Empirical Bayes o PSM para controlar tendencias generales y RTM. 

\section{Evidencia empírica en la Ciudad de México}

\noindent No es la primera vez que se analiza el impacto de las cámaras de velocidad en la Ciudad de México. Dos trabajos destacan por el uso de métodos cuasi–experimentales: Quintero et al.\ (2023), que evalúan el impacto de paquetes de política a nivel alcaldía, y Madrigal y Rodríguez (2021), que estudian efectos locales en el entorno inmediato de los dispositivos.

Quintero et al.\ (2023) analizan dos cambios de política: el programa de \emph{Fotomultas}, introducido en diciembre de 2015, que combinó límites de velocidad más bajos, multas económicas más altas e instalación de radares de velocidad y cámaras para diversas infracciones; y el programa de \emph{Fotoc\'ivicas}, implementado en 2019, del que ya se habló en un inicio y que es objeto de estudio en esta tesis. El diseño empírico es una serie de tiempo interrumpida con regresión binomial negativa, estimada por separado para tres resultados semanales: colisiones totales, colisiones con lesionados y muertes por hechos de tránsito. Para las colisiones utilizan siniestros reportados a AXA México, ponderados por el total de vehículos asegurados a partir de información de la Asociación Mexicana de Instituciones de Seguros (AMIS); para la mortalidad emplean registros de defunciones del Instituto Nacional de Estadística y Geografía (INEGI). Entre los controles explícitos se incluye, por ejemplo, una variable indicadora para las semanas con desabasto de combustibles en 2019, así como términos de Fourier para capturar estacionalidad.\\

Un rasgo metodológico importante de este estudio es la elección de la alcaldía como unidad de tratamiento: si al menos una cámara se ubica en su territorio, toda la alcaldía se clasifica como tratada. Esta decisión facilita la comparación entre municipios con y sin dispositivos, pero diluye los efectos muy locales alrededor de los puntos de control que otros trabajos han documentado. En términos de resultados, para el paquete de 2015 no encuentran cambios estadísticamente significativos en el número total de colisiones ni en las colisiones con lesionados, pero sí una ligera mejora en la mortalidad: la tendencia de defunciones pasa de ser prácticamente plana a mostrar una caída semanal pequeña pero sostenida. Para las reformas de 2019, en cambio, no se observa un cambio en el total de siniestros, pero sí un incremento en la tendencia de colisiones con lesionados (del orden de 1.5\% semanal) y un aumento en la tendencia de mortalidad cercano a 2.7\% semanal. Los autores interpretan estos hallazgos como evidencia de que la eliminación de multas económicas y su sustitución por sanciones cívicas, difíciles de hacer efectivas, redujo la capacidad disuasiva del sistema.\\

Madrigal y Rodríguez (2021) abordan el mismo par de programas pero con una estrategia empírica y una escala espacial distintas. Utilizan datos de incidentes viales confirmados del C5 de la Ciudad de México entre 2014 y 2020 y estiman el impacto de \emph{Fotomultas} y \emph{Fotoc\'ivicas} mediante un diseño de Diferencias en Diferencias. A diferencia de Quintero et al., aquí el tratamiento se define directamente sobre las vialidades donde se instalaron los dispositivos: para cada cámara construyen un área de 300 metros antes y 300 metros después a lo largo de la vía, y consideran como grupo de control los cruces de la red primaria sin cámaras, bajo la lógica de que es en ese tipo de vialidades donde típicamente se colocan estos equipos. Las series se agregan a nivel mensual y se estiman efectos para tres resultados: incidentes totales, con lesionados y con fallecidos, sin covariables adicionales más allá de las tendencias comunes capturadas por el esquema de diferencias en diferencias.\\

Con este diseño más local, los autores encuentran patrones distintos a los de Quintero et al.\ (2023). Para \emph{Fotomultas} no identifican reducción en incidentes totales ni en incidentes con lesionados y, en cambio, estiman un aumento estadísticamente significativo en hechos fatales en las inmediaciones de los dispositivos respecto a los cruces de la red primaria sin cámaras. Para \emph{Fotoc\'ivicas}, por el contrario, documentan reducciones locales en incidentes con lesionados (significativa a un nivel de confianza aproximado de 85\%) y en hechos fatales (alrededor de 90\% de confianza), sin cambios en el total de incidentes. Su interpretación es que el rediseño del programa y la reubicación de cámaras hacia tramos con mayor siniestralidad previa mejoraron el alineamiento entre la ubicación de los dispositivos y el riesgo real, lo que permitió observar efectos localmente favorables aun cuando estos no sean visibles al agregarlos a nivel de toda la ciudad.\\

En conjunto, la evidencia para la Ciudad de México sugiere que los resultados de los programas de fotoinfracciones dependen de al menos tres elementos: (i) el tipo de sanción y la certidumbre de castigo (multas económicas contra sanciones cívicas difíciles de hacer cumplir), (ii) la escala y unidad de análisis utilizada (alcaldía contra entorno inmediato de los dispositivos) y (iii) la forma en que se seleccionan y reubican los sitios de control en función de la siniestralidad previa.

Estos trabajos dejan abierta una pregunta central para la CDMX: si las fotocívicas generaron un efecto causal localizado en el entorno inmediato de los radares una vez que se corrige explícitamente el sesgo de selección en la ubicación de las cámaras y los cambios exógenos en la movilidad. Quintero et al. identifican efectos agregados de política, pero no efectos locales alrededor de dispositivos específicos; Madrigal y Rodríguez se acercan a esa escala local, pero con un contrafactual menos justificado y sin controles explícitos para los cambios de exposición vehicular durante la pandemia. En ese espacio se ubica esta tesis: evaluar el efecto local de las cámaras con una unidad espacial consistente con el mecanismo del tratamiento, un contrafactual construido de forma explícita y controles de movilidad que permitan una mejor interpretación causal.

\section{Diseño e implementación}

\noindent La evidencia revisada coincide en que las cámaras de velocidad sólo son efectivas cuando se insertan en una estrategia más amplia de gestión de la velocidad y de \emph{sistema seguro}. Tanto la \emph{Guía para el control de la velocidad} de la Secretaría de Salud (García Cruz, 2018) como la guía del Banco Mundial y el GRSF (\emph{Global Road Safety Partnership} por sus siglas en inglés) sobre preparación para cámaras y otros sistemas automatizados (Job et al., 2020) parten de la misma premisa: la velocidad es el factor central del riesgo vial, y la probabilidad de muerte crece de manera exponencial con pequeñas variaciones de la velocidad de impacto. Por ejemplo, la probabilidad de que un peatón muera tras un atropellamiento pasa de alrededor de 10\% a 30 km/h a casi 90\% a 50 km/h. Un programa de cámaras que ignore esta relación, o que se diseñe sólo como herramienta recaudatoria, difícilmente generará beneficios sostenidos en seguridad vial.\\

\subsection{Enfoque de sistema seguro}

\noindent La guía de la Secretaría de Salud enfatiza el enfoque de \emph{sistema seguro}: se asume que las personas son vulnerables y cometen errores, por lo que la infraestructura, los vehículos, los límites de velocidad y la legislación deben diseñarse para compensar esos errores y mantener las velocidades de impacto por debajo de umbrales tolerables para el cuerpo humano. Esto supone, entre otros elementos:

\begin{itemize}
    \item Establecer límites de velocidad coherentes con la función de cada tipo de vía, el tipo de usuarios que la utilizan y la calidad de la infraestructura.
    \item Alinear los límites con el diseño vial: cuando el límite es mucho menor que la velocidad sugerida por la sección, la forma y la visibilidad de la vía, se generan flujos muy heterogéneos y conductores que perciben los límites como arbitrarios.
    \item Complementar el marco normativo con medidas de ingeniería y con intervenciones de educación y comunicación. Las campañas aisladas tienen poco efecto.
\end{itemize}

Dentro de este marco, las cámaras de velocidad se conciben como una herramienta para aumentar drásticamente la probabilidad de detección de excesos de velocidad en un entorno donde, sin automatización, el riesgo objetivo de ser sancionado es muy bajo. La Guía destaca dos estudios donde se muestra que en Suecia solo 3 de cada 10,000 incidentes son capturados (Nilsson y Engdal, 1986) y que en Noruega solo 1 de cada 1,000 es detectado (Endersen, 1978).\\

\subsection{Condiciones mínimas para la implementación}

\noindent En cuanto a la implementación de radares, la guía de preparación de Job et al.\ (2020) ofrece un marco simple para evaluar si una ciudad está lista para usar cámaras de velocidad y otros sistemas automáticos. Además, los autores insisten en que no es necesario esperar a tener un sistema perfecto: basta con contar con capacidades básicas para operar, validar mediciones y tramitar infracciones de forma confiable; bajo esas condiciones, los beneficios en seguridad vial tienden a superar los costos de arranque.\\

La guía destaca tres frentes. En lo político y legal, recomienda dejar claro quién responde por la infracción (\emph{owner onus}: si no se identifica al conductor, la sanción recae en el propietario), y construir un marco normativo sencillo y transparente que permita el uso de cámaras, defina estándares mínimos de calidad de los equipos y asegure que el programa se comunique como una política de seguridad vial, no meramente recaudatoria. En la estrategia de control, sugiere combinar cámaras visibles (\emph{overt}) con otras menos visibles (\emph{covert}) y priorizar sitios con historial de choques graves, velocidades típicas altas o riesgo alto de colisiones. Finalmente, subraya el potencial del control de velocidad media por tramo (\emph{average speed enforcement}), que al medir la velocidad a lo largo de un segmento reduce el incentivo a frenar sólo antes del radar y favorece velocidades más homogéneas, sobre todo en vialidades largas y sin accesos intermedios.\\
